\chapter{第五集 \quad 择中}

\section{正文}


黄河在内蒙古高原和黄土高原之间形成的大“几”字形，是造化之力在神州大地书写出的最壮阔一笔。紧邻这个“几”字形大弯的东侧，是吕梁山脉、太行山脉和它们之间的一连串狭长盆地。距今4300年前，在其中的临汾盆地出现了一个曾辉煌一时，却因岁月久远而湮没无闻的古国，它的都邑就在今天山西襄汾陶寺村一带，研究者遂以地为名，称之为陶寺。雾中，考古工作者们忙碌却有序地清理着陶寺遗址的一角。随着覆盖其上的泥土逐渐被揭去，这里赫然显露一场灾难的证据，4000年前的一位年轻贵族，侥幸从这危机重重中逃脱，他宁愿投身那似乎终古不见天日的密林中，也不愿再面对那曾经繁盛的都邑里发生的一切。

对于整个陶寺都邑中的贵族来说，公元前2000年，是动荡与离丧的开端。曾经守卫森严的夯土城墙被破坏、被废弃，最终也未能护住城中的安宁。耗费无数人与物才建就的宫殿已被粗暴的拆毁，昔日宫室中，那戳印精美图案或绘制蓝彩的白灰墙，曾是令造访者炫目的华丽，已成为散落在泥土中的建筑垃圾。政令所出的宫殿，被改造成制造骨器和石器的手工作坊。最令人触目惊心的，是作坊旁的垃圾沟内，那散乱混杂在骨石废料中的人骨。就连早已作古的人们也没能逃过噩运，陶寺的王墓和贵族墓中，大多有直捣中央棺室的扰坑，死者的头颅、碎骨被随意扬弃于墓内各处。种种迹象表明，这非是为求财而来的盗墓，而是一场宣泄式的毁墓虐尸行为。

这场祸及先人和生者的暴乱因何而起？沉默的陶寺遗址并未告诉我们答案。日夜奔流的汾河边上，陶寺之邑曾经强盛一时。宽近八米的城墙维护着280万平方米的土地，守城的卫兵需要快走一小时，才能完成巡城一周的任务。当时的黄河流域和长江流域，一座座城邑拔地而起，风云竞涌，古国林立，俨然正是传说中那万邦并伫，英雄辈出的景象。然而放眼彼时之世间，在规模上能比肩陶寺城址的也并不多。这样庞大的工程规模和城区空间，足能展露出陶寺社会的综合实力，不难想见城内的人口数量之巨，统治者的社会动员力之强。22号墓是陶寺的大墓之一，在经历过发冢之劫后，排场依旧惊人。

墓穴至今仍深7米，陡直的墓壁上装饰着五周平行的手抹草拌泥宽带，或许是对生前居室装修风格的模仿。这壮观的大墓里盛装着的，不只是陶寺贵族曾经的荣光和尊严，还有属于那个黄金时代的迷梦。在曾经的安葬仪式上，人们把一具船形木棺安放在深深的墓底正中。木棺由一根整木挖凿而成，通体红彩，鲜艳夺目，上覆布置棺罩。一个完整的公猪下颌骨，被高高的放置在墓主人头端的墓壁中央。以其为中轴，这面墓壁下方两侧，各倒置着三件彩漆柄的玉石钺和戚，兵器的陈设显现出仪仗的威风。棺木的左侧与墓壁间，排列着4柄青石刀和7块木案板，石刀下的骨骸和朽灰表明，下葬时案板上还盛放着鲜猪肉。

墓主人足端一侧，摆放着10头对半劈开的猪，墓的四壁底部掏有11个壁龛，里面放置精美的玉器、漆器和彩绘陶器。纵然迄今探明的陶寺墓葬已有数千座之多，其中九成都是仅能容身、空无一物的小墓。像22号墓这样规模的大墓实属少数，这反映出陶寺古国的金字塔式等级结构。王族高居塔尖，在其下依次是贵族、平民、赤贫乃至非自由人，社会分化相当严重。权贵们的大墓即便在幽暗的地下，也极尽所能的炫耀着自己对于珍稀资源的掌握，彰显着非比寻常的身份。在大墓中，屡屡可见融汇四方之风的华美器具，玉琮和玉璧宛然可见黄河上游齐家文化的印记。

镂空兽面，散发着长江中游后石家河文化的神秘气息，鼍鼓之上似能听到来自海岱地区礼乐文化的节拍；陶器上的彩绘图案，甚至依稀有东北地区文化之遗风。当然，繁盛如陶寺，又岂会满足于仅以远方殊物表现自己的权威？他们早已掀起新的文化浪潮，建构了自己独特的制度逻辑。在陶寺的几座大墓中，考古工作者发现了一种从未见过的彩绘陶盘，每墓仅有一件，全都统一放置在墓主右侧偏上部。这样的现象，表明陶盘之珍罕，它们或许是只有特定人群才能拥有的，具有特殊含义的仪式用物。这些陶盘通体覆有一层黑色磨光陶衣，盘内以朱红色颜料绘制一条盘曲的动物纹。

此动物，长躯麟身、方首圆目、巨口长舌、无角无爪，似蛇非蛇，似鳄非鳄，显然并非尘世中物，而是陶寺人揉合了多种动物的特点，创造出来的灵兽。它的形态与中国人所熟悉的龙十分相似。在其中一座随葬陶龙盘的大墓内，墓主足端还放置有一组鼍鼓和石磬。鼍，是扬子鳄的古称，鼍鼓即蒙以鳄鱼皮的木鼓，鼓腔呈竖立筒状，高1米，以树干挖制而成，外壁通体施彩绘。石磬是打击乐器，后世的磬大小相递，成套使用，是以又称编磬。陶寺之磬却仅为单件，因此称为“特磬”。另有数座大墓也在墓主足端放置了这样的特磬与鼍鼓或陶鼓之组合，磬音清越，鼓声嘭嘭，应节而起。

陶寺的红铜铃是迄今所知中国历史上的第一件金属乐器，它的出现，预示着“金石之声”时代的即将到来。金声即铜声，铜与玉石和鸣，正是夏商周三代音乐文明的重要内涵。鼍鼓声沉，石磬声清，铜铃声在更高的频率上带来了另一个声部、另一种音色的和鸣。音乐也呼应着世间烂漫生长的万物，众声皆备，这是初萌于陶寺的合乐之声。铜铃虽小，制作工艺却比大件的实心工具和兵器要更复杂，难度也更大。铸造普通的实心铜质工具仅用单范便可以完成，而铜铃的腹腔中空，铸造时必须使用复合范，也就是多于两块以上的范。因此，陶寺的红铜铃是迄今所知年代最早的，以复合陶范铸造的完整铜器，这种铸造技术也正是后来中国青铜时代的标志性工艺。

陶寺文化所处的龙山时代，在时间上正接续在青铜时代之前。陶寺所处的晋西南，在空间上本就是大中原的组成部分，而且恰恰位于中原与北方交流的重要孔道上。虽然陶寺的铜器是使用天然铜铸造的红铜器，有别于后来的青铜合金，但从时空两重因素推测，陶寺一定曾在中原青铜文明的崛起过程中，扮演过重要的角色。四方的资源技术汇集到陶寺，也造就了这座古城规划有度的布局，城内功能分区明确，宫殿、贵族生活区、祭祀场所、手工业作坊、仓储区和不同规模、规格的墓地一应俱全，是这一时期大中原地区，“城市要素”最为完备的巨型都邑。传说与现实的交织引人遐想，莫非尧曾在此驻足？

西倚黄河，东望太行的陶寺，一度曾如此辉煌。这里社会内部等级划分明确，礼乐器群逐渐创备，礼制作为建立新秩序的重要支柱，已经初行于世，耀熠着即将到来的青铜时代。而与兼收并蓄的广阔胸襟和都邑布局的宏大气势，形成鲜明对比的是，考古调查表明，陶寺文化的聚落，基本上只分布于其都邑所在的临汾盆地周边，山河表里，是可凭恃的天险，却也是可屏翳的阻隔。大邑小国的陶寺，数以万计的人口聚集于辐辏之地，至中晚期，不同阶层间又发生着激烈的斗争。距今4000年后，陶寺文化和陶寺都城开始走向衰亡。

石峁，是黄河的“几”字形大湾所包围的黄土高原上生长起来的一个古国。它处在一片被称为“半月形地带”的内弧一侧。半月形地带是中国大陆自东北、西北而西南的一道美丽弧线，形似一枚嵌在东半部的第一、二级阶地外围的新月，月牙朝向孕育古老华夏文明的黄河中下游和长江中下游。这道半月形弧线上，层峦叠嶂，连绵而跌宕。自东北的大兴安岭、内蒙古的阴山、宁夏的贺兰山，直至川西、滇北的横断山。这片半月形地带的存在，最初是由于中原农耕人群向周边的扩散而形成的。因为人口膨胀，许多人需要向外寻找适宜生存的地方。他们通过智慧、双手以及人彼此间的协作，在与平原迥异的生态环境中也发展出了新的生存策略。他们将黄河流域和长江流域的粟黍、稻作农耕技术带到这些山地里春耘秋收，形成星点散落其间的各个定居社会。

而这个地带，也成为了与欧亚大陆其他文化产生碰撞、融合的前沿阵地。当黄河中下游、长江中下游的大部分地区先后进入古国形态时，地处更外围的半月形地带也迎向了时代的风潮。在毛乌素沙漠东南缘，黄土高原北端的黄河西岸，秃尾河畔的黄土梁峁和剥蚀山丘上，海拔千米之处屹立起一座石峁古城。这座始建于4300年前，面积达400万平方米的巨型城址，是那个时代陕北、晋西至蒙中地区的中心。由于在地面以上留存了相当高的一部分，曾一度被后人误认为是战国时期的秦长城附属边塞遗址。万万没想到，一座史前的城址竟能保存的如此完好，秋日的夕阳，只留一抹染红石峁城。

这石城之壮观，为史前东亚地区所罕见，其核心区是高大巍峨的“皇城台”，这是座四周以石块包砌成阶状护坡的台城，顶部十分宽敞，足有8万余平方米。石峁城有内外二重城墙，内城将“皇城台”包围其中，城墙依山势大致呈东北、西南向分布，面积约210万平方米，外城利用内城东南部墙体继续向东南方向扩筑一道弧形石墙，形成封闭空间，城内面积约190万平方米。石峁内外城，各发现四座城门，多数城门地表仍可见墩台、瓮城、角楼等各类附属建筑，甚至还有疑似后世的“马面”结构，这些多是适于布设驻防力量的设施，将石峁城层层卫护，防御森严。

距今4000年前的深秋，夜幕覆于千万道坼裂的黄土梁峁之上，天地间一点点微光徐徐闪烁，火把之间，石峁首领正缓慢地、庄重地登上皇城台，登得越高，就离苍穹越近。明月孤悬，台城高矗，月光、人影与火光摇曳交映，流动过立于地面的神面纹圆形石柱，流动过嵌在石墙中的玉器，以及那些人面、动物和神面的石雕。环眼神面，是源自江汉之间石家河古城的庄严肃杀，人射马，是来自辽阔草原的生机勃发，两条龙盘曲如环，身披鳞甲，充满着莫测的力量。他伸手摩挲着石雕龙身，神龙犹如被唤醒，在暗夜的微光里闪耀着灵动的生命感。这些石雕散落着嵌于石墙间，有的甚至被倒置，但丝毫不影响那粗犷线条凿刻出的原始美感，石墙缝间，还常有故意嵌入的玉制刀钺，扁薄锋利的带刃玉器原是古国时代各区域首领下葬时，显示军权、王权的身份标志。

在石峁，它们却被广泛用于皇城台的营建、贵族宫殿的奠基等公共工程中。藏玉于墙，这是石峁人独特的习俗，玉所具备的透明、坚韧、锋利的特质，仿如一种神力，同石雕图像一起守卫着这座城，给予特别的护佑。首领持玉璋向天祝祷，群山之间，他听见秃尾河汨汨流淌的声音，听见虫鸣和风呼啸的声音，他对这片土地如此熟悉，即便身处光线微弱的黑夜里，也知道月色底下那山谷之间的小型石城——都在何方。往东南，渡过黄河，翻越吕梁山，便是大邑陶寺，再往东往南是更广阔的中原、海岱和江汉。中原，是石峁人的根脉，仰韶时代，祖辈们从太行山以东和关中地区北上移民至此，带来了故土的文化，也在新的家园生长出新的传统。

从石峁往北，是辽阔的欧亚草原，石峁人从那里引进了绵羊、山羊和黄牛，过上了半农半牧的生活。他们还从那里引进了砷铜合金或锡铜合金的青铜刀，获得了坚利的工具，石峁的一些石雕也能在广义阿尔泰地区的石雕形象中找到母本。但4000年前的那个夜晚，让石峁首领陷入沉思的，并不是自何而来的问题，而是往何处去的答案。从对古气候的研究中，我们知道，当时的全球气候正值又一轮周期性变化的阶段，进入了干冷时期。这轮气候变化带来的生存压力，尤其被处于农牧交界地带的石峁集团深刻感知。环境的恶化、人口的增长、资源的紧张，无不是随时可能压垮北方脆弱生态的最后稻草。各聚落之间，更是为了生存而征伐不休。

某个聚落的人家里，孩子因为饥饿而哭泣难以入眠，怀抱孩童的母亲耐心的哄劝，父亲在旁叹了声气，望向窗外的浓稠夜色，这一刻，皇城台上的石峁首领也正望向空荡荡的天地间，忧思重重。在危机时刻，趋利避害，凝聚族群，自是当务之急。但他也决定听从智者们的建议，向外去寻找更多的资源。何去何从，需祈祷上天，敬告祖先！他将口簧含入嘴里，一阵神秘之声响起，这声音震荡着山谷的夜色，仿若以此，达成与祖先神祇的某种神秘共识，南下，是唯一的选择。他相信漫漫长夜之后，石峁城将在黎明中复苏。

龙山文化晚期，征战、对峙、暴力、冲突在中原和北方地区成为常态，活跃于史前中国大地异彩纷呈的区域性古国先后走向衰落，江汉平原和江南地区人烟稀少，山东龙山文化逐渐式微。在陶寺湮灭之后，石峁也开始衰落，中原其他城池纷纷退出历史舞台。而时代的战车，将由新的王者驾驭。距今3700多年前，洛阳盆地北面，邙山的一处高岗之上，夏王少康面南而立，洛河经山下浩浩东去，南面，伊水两岸双峰耸立，如同门阙；西面，晨雾尽头是狭长的函谷，直通关中盆地。根据古文献记载，夏是中国第一个王朝，是王权世袭的开端，开辟九州的禹，为天下苍生栉风沐雨，披荆斩棘。此后的数代王者，有奋发与开拓之主，也有失德和流亡之君。几十年前，先王太康在一次田猎活动中，被东夷的后裔突袭，下落不明。太康失国之后，他的五个兄弟辗转流亡，在洛水边悲吟《五子之歌》。

“明明我祖，万邦之君。有典有则，贻厥子孙。民可近，不可下。民惟邦本，本固邦宁。”如今，东夷之人已经被降服，少康夙夜匪懈，终成复兴大业，他重回洛阳盆地，定都斟鄩。位于中原腹地的河南偃师二里头遗址，历经半个多世纪的田野考古工作，让人们看到了传说中夏都斟鄩的影子。它的现存面积达300万平方米，周边还环绕多个区域中心、次级中心及众多更小聚落，组成了金字塔式的聚落等级结构和众星捧月式的分布格局。这是一个构思宏大、规划缜密、布局严整、等级有序的超级都邑。文献只记录了它的地理方位，考古发现终于展露出它动人的细节。

二里头都城的中心，两条南北向和两条东西向道路纵横交错，呈方正规整的“井”字形，构成中国最早的城市主干道网，在连接交通的同时，又分割出不同的功能区。居中的是大中型夯土建筑基址集中的宫殿区；官营手工业作坊区和祭祀区，分别位于宫殿区的南、北两侧；贵族居住和墓葬区拱卫在宫殿区周围；再向外的部分，则是一般性居住活动区域，择天下之中而立国，择国之中而立宫，择宫之中而立庙。二里头都城所处天下之中的位置，体现出早期国家等级分明的社会结构和秩序井然的统治格局。而九宫格式的都城规划布局，又与《禹贡》以山川划分九州的“天下观”颇为契合，足以显示二里头王者在都邑建设中“辨方正位、体国经野”的政治抱负。

中条山谷中，王朝的矿工们摸索行进，他们的使命是在这群山中找到新的铜矿。那些往来于运城盐湖和二里头王都之间的盐队，常常需要翻越中条山，其中有人曾在连日雨后目睹远处山石上绽开绿花的奇景。矿工们知道，这或许正是铜矿存在的标志，是以扎入深山四处找寻。此前已有入山寻矿的人殒身于此，于是他们格外注意脚下的路。此时山雨突至，只好寻了一处岩厦暂避。连日的疲惫中，众人沉沉睡去，唯有一人在火堆边默默疗伤，一只翠蓝的蝴蝶蓦然飞起，使他心头油然生出一种奇异的异象，便不由自主地起身追逐。一路追寻而来的矿工，看见蝴蝶渐渐地落在对面的山崖上，那处山体因下雨而发生了滑坡，就在他的眼前，那石上赫然绽开一小丛一小丛铜绿之华，这正是他们要为王朝寻找的涉及国家命脉的战略资源。

山西绛县的西吴壁矿冶遗址，仍留存有二里头文化晚期的冶炼遗迹，铜矿原料便来自附近的中条山。这一时期的木炭窑，其功能是用来烧制炼铜使用的木炭，透过炼炉的残迹和炼渣的堆积，似乎就能看到铜矿石在炉中熔化，沉入炉底，化为纯铜。在欧亚草原上，青铜原本仅被用于制造武器、工具、装饰品等小件器物，且多为锻打成型。冶金术传入中原后，迅速发展出以复合范铸造技术为核心的铜容器生产业，新的资源技术与本身的礼制文化系统相结合，创造出独特的青铜文明。二里头遗址发现了迄今所知中国最早的青铜钺，而秉持斧钺之人就是有军事统帅权的首领。

铜爵是一种酒器，也是中国最早出现的青铜容器。青铜是合金，需要将铜与锡等矿石按照一定的比例调配才能冶炼而成。铜矿可在中条山获得，锡矿的主要来源却可能是大兴安岭南麓地区或是江西、湖南等长江以南地区。这意味着如要大规模冶炼青铜，就必须要具备远距离的矿料开采和运输的能力，加之铜爵造型复杂，涉及高精的铸造技术，这一切唯有具备比此前任何一个时代更强大的社会调动和资源控制能力，才有可能做到。二里头都邑的贵族们不仅垄断了青铜礼器的生产，也独占了青铜礼器的使用权。青铜，这种人类创造的物质材料，也被转化成为与社会秩序和等级制度相关联的政治资源。工匠们将日常器具的形体赋予它们。

食器、酒器、乐器、武器，一旦与青铜相关联，这些原本具有实用功能的器型，却也成为了一种超脱于日常的象征符号，具备了礼器、祭器的功能，也与权力身份相挂钩。当以烈火铸就的坚利金属，被塑成庄重敦厚的礼器，文雅莹润的玉石却制成了杀伐征戮的武器，这也是礼制观念中奇妙的二元对立。饮食器皿不再日常，成为青铜重器，干戈坚刃敛去杀气，化为玉石韫藉，这最早王朝的礼制观，自兹历千年而不衰。而那些被寄寓了特殊含义的礼器群，爵、钺、璋等器种，持续使用超过千年，进而成为中国古代社会政治文化的重要符号。除此之外，一种传说中的神兽更是成为流传至今的中国图腾，这便是二里头的龙，它的用工之巨、制作之精，前所罕见。

这条龙全长70厘米，眼、鼻皆嵌白玉，其余部位则由2000多片绿松石拼嵌而成，每片绿松石的厚度大多不足一毫米，长度也不过几毫米，原本应是粘嵌在木、革等材料之上，为了便于粘嵌，每个嵌片都制成楔形。嵌片形状几乎无一重复，虽历三千余年，色泽犹翠绿鲜艳。这是已发掘的二里头贵族墓葬中，位置最接近所在宫室中轴线的一座，墓葬虽只有两平方米，随葬品却有上百件，墓主人是一名30多岁的成年男子，大量海贝置于其颈部，头骨上方发现三件斗笠状白陶器，可能为头饰或冠饰的组件。他怀抱绿松石龙形器，手系铜铃，可能是一名巫师。二里头遗址的另几座墓葬中，墓主胸腹间与铜铃一道放置着镶嵌绿松石的铜牌饰，牌饰正面铸出寿面纹，再以数百枚细小的绿松石片细密镶嵌，以大颗的绿松石珠点睛，赫然勾画出一幅神兽形象，这或许正是龙形器的简化和抽象表现。

碧龙飞升，铜铃和鸣，神灵降临，亡灵升天。一如后来《诗经》中所记，周王祭祀宗庙时，“龙旗阳阳，和铃央央”的场景。神巫一足略跛，随着节奏，且行且舞，足迹画出大大的八字，这样的步伐被称为“禹步”，表现着大禹治水行山之艰难。新砦陶器盖上的龙首、陶寺陶盘上的蛇形蟠龙纹和石峁石雕上的龙蛇形象，在此刻合而为一，成为后世中国人精神世界里永恒的图腾。与之相呼应，东亚文明的复兴地区即黄河和长江流域，开始由多元化的古国文明走向一体化的王朝文明。传说舜帝统治天下时，南方的三苗不愿归顺，禹请求带兵征伐，却遭到了舜的拒绝。舜认为轻易动武乃是执政者德行不够的表现，于是偃武修政，执兵器而舞，三苗乃服。这传承至今的干戚之舞，展示的既是礼节仪式之肃穆，也是武力军容之威慑，不战而胜，这是夏王少康的追求。

在九州团聚的盛会上，自各方远道而来的嘉宾们奉上本地的珍惜物产，表达对华夏的敬意，热带海域的海贝贵重而绮丽，长江下游制作的印纹硬陶盉坚实精致，来自西北高原的是片状有刃的大型玉器，成群的黄牛与绵羊从北方草原被驱赶而来，而华夏用以回馈九州之敬意的，是那些象征王朝威仪的礼制和礼器文化。起源于龙山时代的玉石璋，以二里头都邑作为起点，向长江中上游、珠江三角洲一带传播，甚至到达今天的东南亚地区。盉、爵等二里头风格的陶礼器，北可至燕山以北，南可及长江中上游，西可达甘肃、青海一带。这正是最美好的时节，江山丽日，花草芳香，那田间、桑下，城中，殿前，欣欣向荣，一派蓬勃景象。

而少康也有他的职责，他关于这片国土的理想，“明明我祖，万邦之君。有典有则，贻厥子孙。民可近，不可下。民惟邦本，本固邦宁。”黄河九曲，夭矫如龙，九州共望，国之在中。在何以中国的追问中，二里头文化具有划时代的意义。中国在此时完成了从新石器时代到青铜时代的转变，开启了以铜礼器为核心，政治权力与伦理秩序合一的三代文明。二里头的崛起，是中国这个文化共同体内，第一次出现超强的政治文化核心，其影响所及超越良渚、陶寺和石峁等各个时期的代表，呈现融合集中的“王朝气象”。这一文化核心正处于周人所称的“中国”，影响波及四方，确立了中华文明向心力和凝聚力的基点，从此历经沧海桑田，永不离散。
